\documentclass[a4paper, serif, 10pt]{moderncv}

%% moderncv
% color and style
\moderncvcolor{black}
\moderncvstyle{banking}

% renew moderncv command to adapt for CJK colon
\renewcommand*{\cvitem}[3][.25em]{%
  \ifstrempty{#2}{}{\hintstyle{#2}：}{#3}%
  \par\addvspace{#1}}

\renewcommand*{\cvitemwithcomment}[4][.25em]{%
  \savebox{\cvitemwithcommentmainbox}{\ifstrempty{#2}{}{\hintstyle{#2}：}#3}%
  \setlength{\cvitemwithcommentmainlength}{\widthof{\usebox{\cvitemwithcommentmainbox}}}%
  \setlength{\cvitemwithcommentcommentlength}{\maincolumnwidth-\separatorcolumnwidth-\cvitemwithcommentmainlength}%
  \begin{minipage}[t]{\cvitemwithcommentmainlength}\usebox{\cvitemwithcommentmainbox}\end{minipage}%
  \hfill% fill of \separatorcolumnwidth
  \begin{minipage}[t]{\cvitemwithcommentcommentlength}\raggedleft\small\itshape#4\end{minipage}%
  \par\addvspace{#1}}

%% page layout/margins
\usepackage[top=2.5cm, bottom=2.5cm, left=1.5cm, right=1.5cm]{geometry}

\nopagenumbers{}

%% Babel config for English language
\usepackage[english]{babel}

%% fontspec
\usepackage{fontspec}

\IfFontExistsTF{Linux Libertine}{
  \setmainfont[Ligatures={TeX, Common}, Numbers=Lining, ItalicFont=Linux Libertine]{Linux Libertine}
}{}
\IfFontExistsTF{Linux Libertine O}{
  \setmainfont[Ligatures={TeX, Common}, Numbers=Lining, ItalicFont=Linux Libertine O]{Linux Libertine O}
}{}

%% CTeX
% CJK support, used to show CJK characters in the resume
%
% - heading=false: disable ctex heading
% - punct=kaiming: use kaiming punctuations styles for CJK
% - scheme=plain: use plain scheme, do not override \normalsize font size
% - space=auto: space settings for CJK characters
%
% CJK fonts are resolved per locale language, e.g. Noto Serif/Sans CJK JP
% for Japanese resumes, with fallback to the next region if not installed
%
% ref:
% - http://ctan.mirrorcatalogs.com/language/chinese/ctex/ctex.pdf
\usepackage[UTF8, heading=false, punct=kaiming, scheme=plain, space=auto]{ctex}

\IfFontExistsTF{Noto Serif CJK TC}{
  \setCJKmainfont{Noto Serif CJK TC}
}{}
\IfFontExistsTF{Noto Sans CJK TC}{
  \setCJKsansfont{Noto Sans CJK TC}
}{}

%% line spacing
\usepackage{setspace}
\setstretch{1.125}

%% URL styling - use normal text instead of monospace
\urlstyle{same}

%% Patch \httplink and \httpslink to handle full URLs
% The original moderncv macros always prepend http:// or https:// to URLs.
% This causes issues when the URL already has a protocol (e.g., https://example.com)
% resulting in malformed URLs like https://https://example.com.
% This patch checks if the URL already contains "://" and uses it directly if so.
% Note: We use \str_set:Nx (x-expansion) to expand macros like \@homepage first.
\ExplSyntaxOn
\str_new:N \g_yamlresume_url_str

\RenewDocumentCommand{\httpslink}{O{}m}{
  \str_set:Nx \g_yamlresume_url_str {#2}
  \str_if_in:NnTF \g_yamlresume_url_str {://}
    { \url{#2} }
    { \url{https://#2} }
}

\RenewDocumentCommand{\httplink}{O{}m}{
  \str_set:Nx \g_yamlresume_url_str {#2}
  \str_if_in:NnTF \g_yamlresume_url_str {://}
    { \url{#2} }
    { \url{http://#2} }
}
\ExplSyntaxOff

\name{林彥廷}{}
\title{資深人工智慧工程師 | 大型語言模型與 MLOps 專家}
\phone[mobile]{+886 912 345 678}
\email{yenting.lin.ai@example.com}
\homepage{https://www.yentinglin.dev}

\address{中國臺灣，台北市，台北市，台北市大安區復興南路一段390號，106}{}{}

\extrainfo{{\small \faGithub}\ \href{https://github.com/yentinglin-ai}{@yentinglin-ai} {} {} {} • {} {} {} 
{\small \faLinkedin}\ \href{https://www.linkedin.com/in/yenting-lin-ai}{@yenting-lin-ai} {} {} {} • {} {} {} 
{\small \faMedium}\ \href{https://medium.com/@yentinglin.ai}{@yentinglin.ai}}

\begin{document}

\maketitle

\section{簡介}

\cvline{}{\begin{itemize}
\item 具備 8 年以上人工智慧與機器學習實務經驗，專精大型語言模型、檢索增強生成與 MLOps。
\item 曾領導 5 人跨功能團隊，將 LLM 應用導入企業知識管理、客服與文件摘要流程，提升作業效率 40\%。
\item 熟悉從資料清洗、模型訓練、評估、部署到監控的端到端流程，並具備雲端與地端混合部署經驗。
\item 樂於將技術轉化為可落地的產品價值，擅長跨部門溝通、技術文件撰寫與團隊知識傳承。
\end{itemize}}
\section{教育背景}

\cventry{2015年9月–2017年6月}
        {碩士，資訊工程學系，成績：4.1/4.3}
        {國立臺灣大學}
        {\url{https://www.ntu.edu.tw}}
        {}
        {\begin{itemize}
\item 研究主題聚焦於深度學習與自然語言處理，完成以注意力機制為基礎的文字分類研究。
\item 擔任課程助教，協助機器學習與資料結構課程之教學與作業批改。
\item 與實驗室團隊合作發表 2 篇研討會論文，培養獨立研究與團隊協作能力。
\end{itemize}
\textbf{課程}：機器學習、深度學習、自然語言處理、電腦視覺、資料探勘、演算法設計與分析、分散式系統、大數據分析}

\cventry{2011年9月–2015年6月}
        {學士，資訊工程學系，成績：3.9/4.3}
        {國立清華大學}
        {\url{https://www.nthu.edu.tw}}
        {}
        {\begin{itemize}
\item 奠定紮實的資訊工程基礎，並選修人工智慧與資料探勘相關課程。
\item 參與校內專題競賽，以影像辨識應用獲得佳作。
\end{itemize}
\textbf{課程}：資料結構、演算法、作業系統、資料庫系統、線性代數、機率與統計、計算機網路、程式設計}
\section{工作經歷}

\cventry{2022年1月至今}
        {資深人工智慧工程師}
        {雲智科技股份有限公司}
        {\url{https://www.cloudwise-ai.tw}}
        {}
        {\begin{itemize}
\item 領導 5 人 AI 團隊，規劃與交付企業級大型語言模型應用，服務超過 20 家企業客戶。
\item 設計檢索增強生成（RAG）架構，整合向量資料庫與內部知識庫，將問答準確率提升至 92\%。
\item 建置 MLOps 流程，使用 MLflow、Docker 與 Kubernetes 管理模型版本、部署與監控，縮短上線時間 50\%。
\item 建立模型評估與護欄機制，降低幻覺與敏感資訊外洩風險，並制定內部 AI 使用規範。
\item 與產品、資料工程及資安團隊協作，將需求轉化為可執行技術方案，並培育 3 位初階工程師。
\end{itemize}
\textbf{關鍵字}：大型語言模型、檢索增強生成、MLOps、團隊領導、模型評估}

\cventry{2019年3月–2021年12月}
        {機器學習工程師}
        {睿智數據股份有限公司}
        {\url{https://www.brightdata-ai.tw}}
        {}
        {\begin{itemize}
\item 開發客戶流失預測與推薦系統，透過特徵工程與集成學習，將預測準確率提升 18\%。
\item 使用 PyTorch 與 TensorFlow 建立深度學習模型，並以 Airflow 與 Spark 處理每日百萬筆資料。
\item 將模型封裝為 RESTful API 並部署至 AWS，支援即時推論與批次預測場景。
\item 參與 Agile 開發流程，負責技術文件撰寫、程式碼審查與跨團隊知識分享。
\end{itemize}
\textbf{關鍵字}：機器學習、PyTorch、Spark、AWS、RESTful API}

\cventry{2017年7月–2019年2月}
        {資料科學家}
        {數字智能股份有限公司}
        {\url{https://www.digitintelligence.tw}}
        {}
        {\begin{itemize}
\item 分析行銷與營運數據，建立視覺化儀表板與自動化報表，協助管理層制定決策。
\item 運用統計檢定與 A/B 測試評估行銷活動成效，提出優化建議並帶來 12\% 轉換率成長。
\item 建置資料清理與特徵管線，提升資料品質與分析效率，減少 30\% 手動處理時間。
\item 與業務與產品團隊密切合作，將數據洞察轉化為具體營運策略。
\end{itemize}
\textbf{關鍵字}：資料分析、A/B 測試、SQL、視覺化、統計建模}
\section{語言}

\cvline{中文}{母語或雙語水平 \hfill \textbf{關鍵字}：繁體中文、母語}
\cvline{英語}{完全專業水平 \hfill \textbf{關鍵字}：TOEIC 950、技術文件撰寫、國際團隊協作}
\section{專業技能}

\cvline{機器學習與深度學習}{專家 \hfill \textbf{關鍵字}：Python、PyTorch、TensorFlow、scikit-learn、XGBoost}
\cvline{自然語言處理}{高級 \hfill \textbf{關鍵字}：Transformers、Hugging Face、LangChain、RAG、Fine-tuning}
\cvline{MLOps 與雲端部署}{高級 \hfill \textbf{關鍵字}：Docker、Kubernetes、MLflow、AWS SageMaker、GCP Vertex AI}
\cvline{資料工程}{高級 \hfill \textbf{關鍵字}：SQL、Spark、Airflow、Kafka、BigQuery}
\cvline{程式設計}{專家 \hfill \textbf{關鍵字}：Python、Go、TypeScript、Shell、C++}
\section{榮譽}

\cventry{2023年11月}
        {台灣人工智慧學會}
        {最佳應用論文獎}
        {}
        {}
        {\begin{itemize}
\item 以「檢索增強生成於企業知識問答之應用」獲得最佳應用論文獎。
\item 論文提出混合稀疏與稠密檢索策略，有效提升長文件問答的召回率與準確率。
\end{itemize}}

\cventry{2023年12月}
        {雲智科技股份有限公司}
        {年度傑出工程師}
        {}
        {}
        {\begin{itemize}
\item 表彰在大型語言模型產品化與 MLOps 流程建置上的卓越貢獻。
\item 帶領團隊完成 3 項企業級 AI 專案，並顯著提升客戶滿意度。
\end{itemize}}
\section{證書}

\cventry{2022年3月}
        {Amazon Web Services}
        {AWS Certified Machine Learning – Specialty}
        {\url{https://aws.amazon.com/certification/certified-machine-learning-specialty/}}
        {}
        {}

\cventry{2023年8月}
        {NVIDIA}
        {NVIDIA Certified Associate – AI Infrastructure}
        {\url{https://www.nvidia.com/en-us/training/certification/}}
        {}
        {}

\cventry{2024年1月}
        {Google Cloud}
        {Google Cloud Professional Machine Learning Engineer}
        {\url{https://cloud.google.com/certification/machine-learning-engineer}}
        {}
        {}
\section{出版刊物}

\cventry{2023年11月}
        {以檢索增強生成提升企業知識問答準確率之研究}
        {2023 台灣人工智慧年會論文集}
        {\url{https://www.taiwanai.org.tw/proceedings/2023}}
        {}
        {\begin{itemize}
\item 提出結合 BM25 與向量檢索的混合式 RAG 架構，改善長尾問題的檢索品質。
\item 實驗顯示在企業內部知識庫上，回答準確率較基線模型提升 15\%。
\end{itemize}}

\cventry{2022年6月}
        {輕量化 Transformer 於邊緣裝置之推論加速}
        {中華民國資訊學會通訊}
        {\url{https://www.ics.org.tw/journal/2022}}
        {}
        {\begin{itemize}
\item 探討知識蒸餾與量化技術在邊緣裝置上的效能與精度權衡。
\item 提出適用於資源受限環境的推論優化流程，降低延遲 35\%。
\end{itemize}}
\section{項目}

\cventry{2023年1月–2023年8月}
        {基於檢索增強生成的企業內部知識問答平台}
        {企業知識庫問答系統}
        {\url{https://github.com/yentinglin-ai/enterprise-rag}}
        {}
        {\begin{itemize}
\item 整合 Elasticsearch 與向量資料庫，支援文件上傳、切塊、嵌入與檢索。
\item 設計多輪對話與引用來源機制，提升回答可信度與可追溯性。
\item 以 FastAPI 與 Docker 部署，支援權限控管與使用量監控。
\end{itemize}
\textbf{關鍵字}：RAG、LLM、Elasticsearch、FastAPI}

\cventry{2021年3月–2021年12月}
        {支援多語系與多標籤的客服意圖分類服務}
        {智慧客服意圖分類平台}
        {\url{https://github.com/yentinglin-ai/intent-classifier}}
        {}
        {\begin{itemize}
\item 使用 BERT 與多標籤分類架構，處理客服來文意圖與緊急程度預測。
\item 建立資料標註流程與模型監控指標，持續改善分類表現。
\item 將模型部署為微服務，提供即時 API 與批次推論功能。
\end{itemize}
\textbf{關鍵字}：NLP、BERT、FastAPI、Docker}
\section{興趣愛好}

\cvline{開源社群}{貢獻套件、技術演講、社群經營}
\cvline{戶外活動}{登山、馬拉松、攝影}
\section{志願服務}

\cventry{2021年9月–2023年12月}
        {志工講師}
        {台灣資料科學協會}
        {\url{https://www.twdsa.org}}
        {}
        {\begin{itemize}
\item 定期開設機器學習入門與 Python 資料分析工作坊，累計指導超過 200 位學員。
\item 協助協會規劃年度資料科學年會，負責議程安排與講者邀請。
\item 參與偏鄉程式教育計畫，設計適合高中生的 AI 體驗課程。
\end{itemize}}
    
\end{document}